\documentclass{ximera}[font=12pt]

%Next 2 lines create spaces between items in itemize and enumerate
\setitemize{itemsep=1.5ex}
\setenumerate{itemsep=1.5ex}

%Creates space between paragraphs - Maybe???
\setlength{\parskip}{\medskipamount}


\title{Adding Integers: Financial Reasoning}   % Use xmlatex all -s to compile when SERVE refuses to work.
\author{Kelly Stady}
\license{CC: 0}         % replace with an appropriate license, or set it in xmPreamble

\begin{document}

\begin{abstract}

\end{abstract}

\maketitle
\label{xim:AddIntMoneyActivity}


\section*{Adding Integers Using Financial Reasoning}

Number lines and the Orange-Blue Rule are great for visualizing addition and finding answers to simple problems, but they aren't always practical. For example, we would need a very large number line to find the value of $-185 + 99.$  Instead, we can add integers by thinking in terms of money.


\begin{itemize}
\item \textbf{Assets (+)}:  Money that you \textbf{have}, a gift you receive or a deposit.

\item \textbf{Debts ($-$)}: Money you \textbf{owe}, a bill you pay or a withdrawal.

\end{itemize}


When we add integers that represent assets and debts, we are finding the \textbf{net balance}.
\[ \textrm{\textbf{Net Balance}} = \textrm{\textbf{Assets}} + \textrm{\textbf{Debts}} \]


\textbf{Note:} A statement such as ``a debt of \$50'' is represented mathematically as $-50$. The negative sign is omitted in the phrase because the word \emph{debt} already implies a negative value.


\bigskip

\textbf{Examples:}


\begin{enumerate}
    \item \textbf{Combining Assets:} A child has \$100 ($+100$) in their piggy bank and receives \$30 ($+30$) for doing chores, their net balance is \textbf{\$130}.
    \[ 100 + 30 = 130 \]

    \item \textbf{Combining Debts:} If you have a credit card debt of \$500 ($-500$) and make another credit card purchase of \$100 ($-100$), your net balance is a \textbf{debt of \$600 ($-600$)}.
    \[ -500 + (-100) = -600 \]

    \item \textbf{Partial Cancellation:} Suppose you have \$20 ($+20$) and owe a friend \$25 ($-25$). Your \$20 isn't enough to cover the entire debt, it only cancels \textbf{part} of it.  Your net balance is a \textbf{debt of \$5 ($-5$)}.
    \[ 20 + (-25) = -5 \]

    \item \textbf{Full Cancellation:} If you owe the IRS \$300 ($-300$) and have \$800 ($+800$) in the bank, you can pay the entire debt and keep the change. Your net balance is \textbf{\$500 ($+500$)}.
    \[ -300 + 800 = 500 \]
\end{enumerate}


\bigskip

\begin{quote}
\textbf{Key Observations}

%Take another look at the examples above and see if you notice a relationship between the signs of the numbers that were added and the operation performed.

Take a close look at the examples above. Notice the relationship between the signs of the numbers that were added and the operation we actually performed.

\begin{itemize}
    \item \textbf{Same Signs:} When we add two assets or two debts, we \textbf{add} their values.  The sign of the answer is the same as the sign of the numbers.

    \item \textbf{Different Signs:} When we add an asset and a debt, all or part of the debt will be removed or cancelled, so even though we are technically \emph{adding} integers, the arithmetic performed is \textbf{subtraction}.  The sign of the answer depends on which is greater: the asset or the debt.
\end{itemize}


\end{quote}



\begin{problem}
Your checking account is overdrawn by \$50 and the bank charges you a \$15 overdraft fee.  Write the addition problem to find your net balance.

\begin{multipleChoice}
\choice{$50 + (-15)$}
\choice{$-50 + 15$}
\choice{$50+15$}
\choice[correct]{$-50 + (-15)$}
\end{multipleChoice}


\begin{problem} % Nested Problem that should not be visible until student answers the first part.
Find your net balance.  \textbf{Do not use a calculator.}

\[ -50 + (-15) \begin{prompt} = \answer{-65} \end{prompt} \]

\end{problem}

\end{problem}


\begin{problem}
A soccer team spent \$25 on supplies for a car wash.  By the end of the day, they earned \$175. Write the addition problem to find the team's net balance.

\begin{multipleChoice}
\choice{$25 + 175$}
\choice[correct]{$-25 + 175$}
\choice{$-25+(-175)$}
\choice{$25 + (-175)$}
\end{multipleChoice}


\begin{problem} % Nested Problem that should not be visible until student answers the first part.
Find the team's net balance.  \textbf{Do not use a calculator.}

\[ -25 + 175 \begin{prompt} = \answer{150} \end{prompt} \]

\end{problem}

\end{problem}


\begin{problem}
It's your birthday!  You receive \$300 as a gift, but you owe your roommate \$500 for rent. Write the addition problem to find your net balance after you use your gift money to pay toward the rent.

\begin{multipleChoice}
\choice{$-300 + (-500)$}
\choice{$300 + 500$}
\choice[correct]{$300+(-500)$}
\choice{$-300 + 500$}
\end{multipleChoice}


\begin{problem} % Nested Problem that should not be visible until student answers the first part.
Find your net balance.  \textbf{Do not use a calculator.}

\[ 300 + (-500) \begin{prompt} = \answer{-200} \end{prompt} \]

\end{problem}

\end{problem}

\begin{problem}
Use the concepts of assets and debts to add.  \textbf{Do not use a calculator.}

\[ 80 + (-2)  \begin{prompt}=\answer{78}\end{prompt} \]

\[ -9 + (-1)  \begin{prompt}=\answer{-10}\end{prompt} \]

\[ -30 + 27 \begin{prompt}=\answer{-3}\end{prompt} \]

\end{problem}


\subsection*{Adding Integers with Different Signs}

When adding integers with different signs, we can think of the process as a
\textbf{battle between positive and negative numbers}, where the winner is the sign of
the answer.

\begin{enumerate}
\item[] \textbf{Battle 1: The answer is positive, so the positive team win!}


\[ 30 + (-20) = 10 \]

\emph{Reasoning:}  If you have \$30 and spend \$20 to see a movie, your net balance is \$10.



\item[] \textbf{Battle 2: The answer is negative, so the negative team win!}


\[ 30 + (-100) = -70 \]

\emph{Reasoning:} If you have \$30 and a credit card debt of \$100, your net balance is a debt of \$70.
\end{enumerate}

Why did the positive team win the first battle and lose the second battle?

%The debt in the second battle was greater than the asset, so the positives lost that one.


In any number battle, the winner is always the team with the \textbf{greatest strength}, regardless of its sign.  Mathematically, we find the \textbf{absolute value} of a number to determine its \textbf{strength} or \textbf{size}.  Recall that the absolute value of a number is its distance from 0 on a number line.  Because it represents distance, it is \textbf{always positive}.


\textbf{Battle 1}
\begin{itemize}
    \item The strength of the \$30 asset is $|30| = 30$.
    \item The strength of the \$20 debt is $|-20| = 20$.
\end{itemize}


Since $30$ is greater than $20,$ the \textbf{positive} team was stronger and \textbf{won} the battle.


\textbf{Battle 2}
\begin{itemize}
    \item The strength of the \$30 asset is $|30| = 30$.
    \item The strength of the \$100 debt is $|-100| = 100$.
\end{itemize}

Since $30$ is smaller than $100,$ the \textbf{negative} team was stronger and \textbf{won} the battle.



\end{document}